\begin{explanationbox}
	\textbf{\textcolor{CExplanation}{一句话直觉}}
	塞瓦定理把“三线共点”的几何条件转化为三条有向线段比的乘积为 $1$，使共点问题可代数计算。

	\medskip
	\textbf{\textcolor{CExplanation}{核心拆解}}
	\begin{description}[style=nextline, leftmargin=2.8em, labelsep=0.5em, itemsep=0.45em]
		\item[\textbf{要点一（共点转比例）：}]
			三线 $AD,BE,CF$ 共点时，利用面积比推出各边分点间的比例关系.
		\item[\textbf{要点二（有向比符号）：}]
			比值采用有向线段，自动处理点在延长线上的符号问题.
	\end{description}

	\medskip
	\textbf{\textcolor{CExplanation}{几何本质（面积比桥接）}}
	三角形中等高三角形的有向面积比等于有向底边比：$S_{\triangle ABD}:S_{\triangle ACD}=BD:DC$，依此类推；三线共点导致三个面积比乘积等于 $1$.

	\medskip
	\textbf{\textcolor{CExplanation}{代数意义（乘积为定值）}}
	不论分点如何，只要三线共点，有向比例乘积始终为 $1$；反之，若乘积为 $1$ 且有两线相交，则三线共点.

	\medskip
	\textbf{\textcolor{CExplanation}{考点价值（解题实用）}}
	\begin{itemize}[leftmargin=*, itemsep=0.5em, topsep=0.4em]
		\item \textbf{考法一（共点性证明）：} 给出线段比，验证乘积是否为 $1$，判定三线是否共点.
		\item \textbf{考法二（未知比例计算）：} 已知两组比，利用乘积 $1$ 求解第三组比.
	\end{itemize}

	\medskip
	\textbf{\textcolor{CExplanation}{顿悟点}}
	\textbf{塞瓦定理本质是“面积比链”的积消，将几何共点转化为代数等式，是数形结合的典范.}

	\medskip
	\textbf{\textcolor{CExplanation}{使用场景}}
	\begin{itemize}[leftmargin=*, itemsep=0.5em, topsep=0.4em]
		\item \textbf{场景一（内分点共点）：} $D,E,F$ 都在三角形边上，此时三线必交于一点.
		\item \textbf{场景二（延长线上点）：} 点位于延长线上，仍需用有向比公式，乘积仍为 $1$.
	\end{itemize}

\end{explanationbox}
